\documentclass[article]{jss}

\usepackage{amsmath,booktabs}
\usepackage[T1]{fontenc}
\usepackage[utf8]{inputenc}

\DeclareMathOperator{\Imag}{Im}

\author{Simon A. Broda\\Institute of Financial Services IFZ\\Lucerne University of Applied Sciences}
\Plainauthor{Simon A. Broda}

\title{Fast Tails of Quadratic Forms in \pkg{QuadraticFormsMGHyp.jl}}
\Plaintitle{Fast Tails of Quadratic Forms in QuadraticFormsMGHyp.jl}
\Shorttitle{Fast Tails in \pkg{QuadraticFormsMGHyp.jl}}

\Abstract{
\pkg{QuadraticFormsMGHyp.jl} evaluates the survival function and the expected shortfall of a quadratic form in a multivariate generalized hyperbolic vector.
The expressions are the Gil-Pelaez integrals derived by \citet{bz:21}.
Version 1.5 of the package replaces an adaptive quadrature, which recomputed a complex Bessel function at every sample of every threshold, with a mapped Gauss--Legendre rule.
The nodes do not depend on the threshold.
The normal-inverse-Gaussian case and the boundary case $\psi=0$ use closed forms, and a list of thresholds is reduced to a Chebyshev series whose degree grows until the last coefficient is negligible.
On a ten-dimensional normal-inverse-Gaussian form with 1{,}401 thresholds, one thread of version 1.4 takes 1.426 seconds after warmup.
The same call in version 1.5 takes 0.124 milliseconds.
A C implementation of the same node rule takes 0.076 milliseconds.
}

\Keywords{generalized hyperbolic distribution, quadratic form, expected shortfall, Gauss--Legendre quadrature, \proglang{Julia}}
\Plainkeywords{generalized hyperbolic distribution, quadratic form, expected shortfall, Gauss-Legendre quadrature, Julia}

\Address{
  Simon A. Broda\\
  Institute of Financial Services IFZ\\
  Lucerne University of Applied Sciences\\
  E-mail: \email{simon.broda@hslu.ch}\\
}

\begin{document}

\section{Introduction} \label{sec:intro}

Let $L=a_0+a^\top X+X^\top AX$, where $X$ is a $d$-dimensional generalized hyperbolic vector.
The law of $X$ has the representation $X=\mu+W\gamma+\sqrt{W}\,CZ$ with $Z$ standard normal and $W$ generalized inverse Gaussian with parameters $(\lambda,\chi,\psi)$.
Tail probabilities and partial moments of $L$ are the inputs of expected-shortfall calculations for portfolios whose return is locally quadratic, and of the sampling distribution of the two-stage least squares estimator under multivariate $t$ errors.
\citet{bz:21} give both quantities as Gil-Pelaez integrals \citep{gp:51} and a saddlepoint approximation.
The integrals are one-dimensional, but the integrand calls the modified Bessel function of the second kind at a complex argument.

\pkg{QuadraticFormsMGHyp.jl} is a \proglang{Julia} \citep{julia} implementation of that calculation.
Through version 1.4 the exact branch asked \pkg{QuadGK.jl} for two improper integrals at every threshold, and each integrand evaluation called \code{besselk}.
On the ten-dimensional normal-inverse-Gaussian example of Section \ref{sec:time}, one thread takes 1.426 seconds to return 1{,}401 tails.
The same numbers come out of a fixed 32-point rule in a fraction of a millisecond, because the expensive part of the integrand does not depend on the threshold and, for this parameter value, the Bessel factor is elementary.
A \proglang{C} library already evaluates that rule.
The \proglang{Julia} code is compiled as well, so after the first call it can do the same work.
Version 1.5 is that port.
The public function is unchanged: \code{qfmgh} still returns the survival probability $\Pr(L>x)$ and the conditional expected shortfall $\mathrm{E}[L\mid L>x]$.

\section{The integral} \label{sec:integral}

Write $\kappa$ for the constant term that remains after the spectral reduction of \citet{bz:21}, and write $q=x-\kappa$.
The survival probability and the expected shortfall are
\begin{equation}
  \label{eq:tails}
  \Pr(L>x)=\frac12+\frac1\pi I_c(q),
  \qquad
  \mathrm{E}[L\mid L>x]
  =\frac{\tfrac12 M_2(0)+\pi^{-1}I_p(q)}{\Pr(L>x)}+\kappa.
\end{equation}
Here $M_2(0)=\mathrm{E}[L]-\kappa$ is known in closed form from the moments of $W$, and $I_c$ and $I_p$ are imaginary parts of integrals along the positive frequency axis.
The integrands are of the form $u^{-1}\Imag\exp(\ell_\lambda(\chi(u,q),\psi(u)))$, where $\ell_\lambda=\log k_\lambda$ is the log of the generalized-inverse-Gaussian normaliser, analytically continued.
The characteristic function has already been reduced to a sum over the eigenvalues of the symmetrized quadratic form.
Those eigenvalues, and every sum that does not involve $q$, are computed once.

The factor $u^{-1}$ makes a naive truncation delicate, and the integrand oscillates.
An adaptive Gauss--Kronrod rule handles both, at the cost of a new Bessel evaluation at every accepted sample.
For a normal-inverse-Gaussian law that sample is an elementary function, and for every other law it is a complex \code{besselk}.
Either way the adaptive rule repeats the work at every threshold.

\section{A threshold-independent rule} \label{sec:rule}

\subsection{Mapped Gauss--Legendre nodes}

The change of variables $u=(v/(1-v))^2$ sends $v\in(0,1)$ onto $u\in(0,\infty)$ and removes the endpoint singularity.
The image of a comparison tail is cut at an upper limit $u_b<1$ chosen so that a bound on the omitted integral is at most $10^{-10}$, and $u_b$ is clamped into $[10^{-6},1-10^{-6}]$.
Gauss--Legendre nodes on $[-1,1]$ are then mapped to $v\in(0,u_b)$.
The quadrature weights absorb the Jacobian and the factor $u^{-1}$, so the sum is a plain weighted sum of the imaginary part.

The number of nodes depends on the law, because the integrand is smoother for some parameter values than for others.
A normal-inverse-Gaussian law ($\lambda=-1/2$, $\chi>0$, $\psi>0$) and a generic law use 32 nodes.
A half-integer order uses 48.
The boundary $\psi=0$, when the skewness channel does not move the second GIG argument off zero, uses 64.
Zero eigenvalues, those smaller than $10^{-12}$ times the scale of the spectrum, are folded into the sums analytically and do not add nodes.
The truncation limit is computed from the full spectrum, before that folding.

\subsection{Closed forms}

For a normal-inverse-Gaussian law the Bessel factor is $K_{\pm 1/2}$, which is a square root and an exponential.
The three orders $\lambda$, $\lambda+1$ and $\lambda+2$ that enter $I_c$ and $I_p$ therefore share one square root per node.
The implementation follows the stable form used by the \proglang{C} code: $\log K_\nu(z)=\log(e^{z}K_\nu(z))-z$, with the closed form substituted for the scaled Bessel value.
On the boundary $\psi=0$ the same three orders reduce to logarithms and gamma values, again one logarithm of the complex $\chi$-argument per node.

Half-integer and generic orders call the scaled modified Bessel function and subtract $z$.
This is the same analytic expression as the series used in the \proglang{C} library.
It is not the same arithmetic, and Section \ref{sec:time} shows the difference in runtime.
The values agree with an adaptive integral to the tolerances in Section \ref{sec:acc}.

One cancellation is worth making explicit.
The anchor $M_2(0)$ multiplies $\mathrm{E}[W^2]$ by the skewness coefficient of the quadratic form.
For an integer $\lambda\le -1$ and $\psi=0$ that moment is infinite.
If the coefficient is exactly zero the product is zero, and forming $\exp(\ell_{\lambda+2})\cdot 0$ produces a \texttt{NaN}.
Version 1.5 drops the product before the exponential when the coefficient vanishes.
A single threshold of the \proglang{C} library returns a finite survival probability and a \texttt{NaN} shortfall on that input.
The \proglang{Julia} shortfall stays finite.
The two survival probabilities agree.

\subsection{Long grids} \label{sec:cheb}

The map from the threshold to $(I_c,I_p)$ is analytic on intervals that do not cross a sharp bend of the tail.
A Chebyshev expansion on the requested range then replaces one integral per threshold by one integral per node of the expansion.
Grids of length 24 or less are integrated directly.
A longer grid starts at degree 20 in the normal-inverse-Gaussian case and at degree 48 otherwise.
The degree increases until the last coefficient of the survival series is at most $10^{-9}$ and the last coefficient of the shortfall series is at most $10^{-9}$ times the scale of the shortfall, or until the degree reaches the length of the grid or 96.
If the last coefficient has not fallen by then, the requested thresholds are integrated directly.
Repeated thresholds, equal to machine precision, are evaluated once.

On the normal-inverse-Gaussian benchmark the series stops at degree 20.
Across 41 thresholds spaced along the grid from 3.5 to 17.5, the interpolant differs from the 32-point rule at that same threshold by at most $3.5\times 10^{-10}$ in the survival probability and $8.2\times 10^{-10}$ in the shortfall.
A fixed expansion of degree 6, which is enough for a very smooth tail and too short for a generic one, misses the direct rule by about $0.02$ on $\lambda=-1.3$ over $[1,20]$.
The adaptive degree is there so that a long grid does not silently return that interpolant.
A $\psi=0$ tail with a corner near the centre of a short interval does not settle by degree 96.
The package then evaluates every threshold.
That call is still a few tenths of a millisecond, because each node is a logarithm rather than an adaptive Bessel sample.

\section[The package]{The \pkg{QuadraticFormsMGHyp.jl} package} \label{sec:soft}

The package is registered and installs with \code{using Pkg; Pkg.add("QuadraticFormsMGHyp")}.
It has no binary dependency.
The exact branch uses \pkg{LinearAlgebra.jl} for the symmetric eigendecomposition and \pkg{SpecialFunctions.jl} for the Bessel and gamma values.
The saddlepoint branch still uses \pkg{ForwardDiff.jl} and \pkg{Roots.jl}, and it is selected with \code{do_spa=true}.
The examples below leave that branch off.

\begin{Code}
using LinearAlgebra, QuadraticFormsMGHyp

x = range(3.5, 17.5; length = 1401)
ccdf, es = qfmgh(x, 0.0, zeros(10), 0.5I(10), Matrix{Float64}(I, 10, 10),
    zeros(10), zeros(10), -0.5, 1.0, 1.0)
\end{Code}

The first argument may be a scalar.
A scalar call returns two scalars.
The matrices are the quadratic-form matrix $A$ and the scale $C$ in the stochastic representation.
The last three arguments are $\lambda$, $\chi$ and $\psi$.
The pair \code{(ccdf, es)} is $(\Pr(L>x),\,\mathrm{E}[L\mid L>x])$.
The upper partial moment $\mathrm{E}[L\,1_{\{L>x\}}]$ is the product of the two.

Half-integer and generic orders split the node loop across \code{Threads.@threads} when the grid has at least four points and \proglang{Julia} was started with more than one thread.
The normal-inverse-Gaussian and $\psi=0$ loops stay on one thread.
Each node is a few arithmetic operations, and splitting them costs more than it saves.
The timings in Section \ref{sec:time} include construction of the eigendecomposition, which is part of every call to \code{qfmgh}.

\section{Accuracy} \label{sec:acc}

Two comparisons are reported.
The first is against \pkg{QuadGK.jl} with relative tolerance $10^{-10}$ on the survival integral and $10^{-8}$ on the shortfall integral, using the characteristic function already implemented for the saddlepoint branch.
The thresholds are single points, so the Chebyshev grid is not involved.
Table \ref{tab:gk} gives the largest absolute gap on each trio of thresholds.
The quadratic form is $L=X^\top(0.5I)X$ in dimension 10, with $\mu=\gamma=0$ and $C=I$.

\begin{table}[ht]
\centering
\begin{tabular}{lccc}
\toprule
Law & Thresholds & Survival & Shortfall \\
\midrule
$\lambda=-1/2$, $\chi=\psi=1$ & $3.5,\,8,\,17.5$ & $3.2\times 10^{-9}$ & $2.1\times 10^{-6}$ \\
$\lambda=-3/2$, $\chi=2$, $\psi=3/2$ & $1,\,8,\,20$ & $2.4\times 10^{-13}$ & $1.0\times 10^{-9}$ \\
$\lambda=-1.3$, $\chi=2$, $\psi=1.1$ & $1,\,8,\,20$ & $4.0\times 10^{-9}$ & $7.8\times 10^{-6}$ \\
\bottomrule
\end{tabular}
\caption{Largest absolute difference between the mapped rule and an adaptive integral.}
\label{tab:gk}
\end{table}

The second comparison is against the \proglang{C} implementation of the same node rule, called through the \pkg{QuadraticFormsMGHyp} extension for \proglang{Python}.
With both implementations on a Chebyshev grid of degree 16, the 1{,}401 normal-inverse-Gaussian survival values differ by at most $6\times 10^{-16}$.
Version 1.5 stops at degree 20 rather than 16.
Its test suite therefore compares a long grid with the direct rule at the endpoints and at the middle of the grid, and requires absolute agreement to $10^{-8}$.
A single normal-inverse-Gaussian threshold agrees with the \proglang{C} value to a unit in the last place.
The same holds for a half-integer order, a generic order, an integer order with $\psi>0$, and the diagonal option portfolio used by \citet{bz:21}, all evaluated at one threshold.

A fifty-dimensional two-stage least squares form, the reduction in which the instrument is a multiple of the first unit vector, is a harder eigenproblem.
Over the $122$ pairs $(\nu,b)$ with $\nu\in\{3,9\}$ and $b$ on a 61-point grid from $-3$ to $3$, the largest absolute gap is $2.0\times 10^{-7}$, at $\nu=3$ and $b=3$, where the survival probability is $0.0723$.
The centre of each grid agrees to $10^{-16}$.
The gap tracks $|b|$ and is consistent with the difference between the two symmetric eigensolvers.

\section{Runtime} \label{sec:time}

Table \ref{tab:time} reports the minimum runtime of a call that both builds the quadratic form and evaluates it.
Each entry is the minimum of repeated calls after a discarded warmup, so the just-in-time compilation is not included.
The machine is an Apple M4 with 10 logical CPUs and 16 GB of RAM, running macOS 26.5.
\code{-t auto} on this machine starts four \proglang{Julia} threads.
The column headed ``10 threads'' was started with \code{-t 10}.
The \proglang{C} times are from \proglang{Python} 3.9.6 calling the extension and include construction of the \code{QuadraticForm} object.
The \proglang{Julia} column headed 1.4 is version 1.4.0 of this package, which is the adaptive integrator.
Version 1.5 is the code described here, on \proglang{Julia} 1.11.7.

\begin{table}[ht]
\centering
\footnotesize
\begin{tabular}{lrrrrr}
\toprule
& \multicolumn{2}{c}{Version 1.4} & \multicolumn{2}{c}{Version 1.5} & \\
\cmidrule(lr){2-3}\cmidrule(lr){4-5}
Problem & 1 thread & 4 threads & 1 thread & 10 threads & \proglang{C} \\
\midrule
NIG, 1{,}401 thresholds & 1426 & 503 & 0.124 & 0.122 & 0.076 \\
NIG, one threshold & 0.729 & 0.732 & 0.033 & 0.033 & 0.040 \\
Half-integer, 200 thresholds &  &  & 0.816 & 0.298 & 0.126 \\
Generic order, 200 thresholds &  &  & 1.017 & 0.288 & 0.087 \\
\bottomrule
\end{tabular}
\caption{Milliseconds per call after warmup.
The normal-inverse-Gaussian rows are dimension 10 with $\lambda=-1/2$ and $\chi=\psi=1$, on thresholds from 3.5 to 17.5.
The other two rows use the same dimension on thresholds from 1 to 20, with $(\lambda,\chi,\psi)=(-3/2,2,3/2)$ and $(-1.3,2,1.1)$.
Blank cells were not timed.}
\label{tab:time}
\end{table}

One thread of version 1.4 spends 1.426 seconds on the long normal-inverse-Gaussian grid.
Version 1.5 spends 0.124 milliseconds, a factor of $1.15\times 10^{4}$.
Four threads bring version 1.4 down to 0.503 seconds.
The normal-inverse-Gaussian path in version 1.5 does not use those threads, and the timing does not move.
The \proglang{C} library takes 0.076 milliseconds on the same grid, so the \proglang{Julia} call is slower by a factor of about 1.6.
On a single threshold the order reverses: 0.033 milliseconds in \proglang{Julia} and 0.040 milliseconds in \proglang{C}.
The remaining gap on the long grid is allocation around a 20-point expansion, not the quadrature.

The half-integer and generic rows are where the \proglang{C} series is faster than \code{besselk}.
Ten threads bring the \proglang{Julia} generic call from 1.017 milliseconds to 0.288 milliseconds.
The \proglang{C} call takes 0.087 milliseconds.
Both are far from the adaptive integrator: a single normal-inverse-Gaussian threshold already costs version 1.4 0.73 milliseconds, and a generic order evaluates several Bessel functions per node.

The same pattern holds for the diagonal option portfolio of \citet{bz:21}, which is normal-inverse-Gaussian.
Version 1.5 takes 0.137 milliseconds for the 1{,}401 thresholds from 3.5 to 17.5, and the \proglang{C} library takes 0.076 milliseconds.
A ten-dimensional Student-$t$ form with $\lambda=-5/2$, $\chi=5$ and $\psi=0$, on 1{,}401 thresholds from 2 to 30, takes 0.404 milliseconds.
The Chebyshev series settles, so this is not the direct fallback.

\section{Summary} \label{sec:summary}

The tail probability and the expected shortfall of a quadratic form in a multivariate generalized hyperbolic vector are one-dimensional integrals.
Once the eigenvalues are known, a mapped Gauss--Legendre rule with a few dozen nodes evaluates both, and a Chebyshev series extends that evaluation to a long grid of thresholds.
\pkg{QuadraticFormsMGHyp.jl} 1.5 is that calculation in \proglang{Julia}.
After warmup it is within a factor of two of a \proglang{C} implementation of the same rule on the normal-inverse-Gaussian benchmark, and four orders of magnitude faster than the adaptive integrator it replaces.
The saddlepoint approximation of \citet{bz:21} is still available for the calls that ask for it.

\section*{Computational details}

Timings were taken on an Apple M4 with 10 logical CPUs and 16 GB of RAM, under macOS 26.5.
\pkg{QuadraticFormsMGHyp.jl} 1.4.0 and 1.5.0 ran on \proglang{Julia} 1.11.7.
The \proglang{C} comparison used the \pkg{QuadraticFormsMGHyp} extension for \proglang{Python}, imported from \proglang{Python} 3.9.6.
Each reported time is the minimum of repeated calls after one discarded warmup.
Version 1.4 was timed twice on the long grid and five times on the single threshold.
Version 1.5 was timed 200 times on the normal-inverse-Gaussian rows and 50 times on the others.
The \proglang{Python} rows are the minimum of 30 calls.
No time includes compilation or process startup.

\bibliography{qfmgh}

\end{document}
